\documentclass{article}
\usepackage{amsmath,amssymb}
\usepackage{xurl}
\usepackage[colorlinks=true,linkcolor=blue,citecolor=blue,urlcolor=blue]{hyperref}
% Shared commands for notes in docs/paper-gaps/.

\DeclareUnicodeCharacter{2080}{\ifmmode{}_{0}\else\textsubscript{0}\fi}
\DeclareUnicodeCharacter{2081}{\ifmmode{}_{1}\else\textsubscript{1}\fi}
\DeclareUnicodeCharacter{2082}{\ifmmode{}_{2}\else\textsubscript{2}\fi}
\DeclareUnicodeCharacter{2099}{\ifmmode{}_{n}\else\textsubscript{n}\fi}

\newcommand{\Lean}{\textsc{Lean}}
\ExplSyntaxOn
\NewDocumentCommand{\leanid}{m}{
  \ifmmode
    \textsf{\detokenize{#1}}
  \else
    \group_begin:
    \urlstyle{sf}
    \tl_set:Nn \l_tmpa_tl {#1}
    \tl_replace_all:Nnn \l_tmpa_tl {\_} {_}
    \exp_args:NV \path \l_tmpa_tl
    \group_end:
  \fi
}
\ExplSyntaxOff
\providecommand{\tr}{\operatorname{tr}}
\newcommand{\tnleanIssueBase}{https://github.com/LionSR/TNLean/issues/}
\newcommand{\tnleanPrBase}{https://github.com/LionSR/TNLean/pull/}
\newcommand{\tnleanDocBase}{https://sirui-lu.com/TNLean/docs/}
\newcommand{\ghissue}[1]{\href{\tnleanIssueBase#1}{GitHub issue \##1}}
\newcommand{\ghpr}[1]{\href{\tnleanPrBase#1}{PR \##1}}
\newcommand{\doclink}[1]{\href{\tnleanDocBase#1.html}{\path{#1}}}

% Dirac notation and the identity operator, as in blueprint/src/macros/common.tex.
% \providecommand keeps note-local definitions authoritative.
\usepackage{dsfont}
\providecommand{\ket}[1]{|#1\rangle}
\providecommand{\bra}[1]{\langle #1|}
\providecommand{\braket}[2]{\langle #1 | #2 \rangle}
\providecommand{\ketbra}[2]{|#1\rangle\!\langle #2|}
\providecommand{\Id}{\mathds{1}}
\providecommand{\one}{\mathds{1}}
\providecommand{\C}{\mathbb{C}}
\providecommand{\MN}[1]{M_{#1}(\C)}
\providecommand{\MD}{M_D(\C)}

% Verdict marker: \gapnote{<kind>}{<status>}. Kind is the policy
% classification (clarification, local-correction, scope-restriction,
% unfaithful, false-source, open-gap); status is open, wip (elimination
% underway), resolved, or historical. Machine-read by the site index;
% prints a verdict line.
\newcommand{\gapnote}[2]{\par\noindent\textbf{Verdict:} #1 (#2).\par}

\setlength{\emergencystretch}{2em}
\title{The Exponent of Hastings' Gaussian Time Average}
\author{}
\date{October 22, 2026}
\begin{document}
\maketitle
\gapnote{local-correction}{resolved}

\section{The source statement}

In the appendix of~\cite{Hastings2007AreaLaw} the Hamiltonian $H$ has a unique
ground state $\Psi_0$, ground energy zero and gap $\Delta E$, and
$\xi=\max(2v/\Delta E,\xi_C)$.\footnote{\path{References/0705.2024/law16.tex},
    lines~126--142 and~766--767.}
With $q=(l/3)\Delta E/(2v)$, the operator
\begin{align}
    P_q=\frac{\Delta E}{\sqrt{2\pi q}}\int e^{iHt}e^{-(t\Delta E)^2/2q}\,dt
    \label{eq:hastings07_gaussian_average}
\end{align}
is said to satisfy
\begin{align}
    \|P_q-P_0\| & \le O\bigl(e^{-l/3\xi}\bigr),
    \label{eq:hastings07_gaussian_average_source}
\end{align}
where $P_0$ is the ground-state projector.\footnote{\path{References/0705.2024/law16.tex},
    lines~792 and~843--848.}

\section{The point at issue}

The operator (\ref{eq:hastings07_gaussian_average}) is diagonal in an eigenbasis
of $H$, with eigenvalue $e^{-qE^2/2\Delta E^2}$ at energy $E$. It is $1$ on the
ground state, and on the excited states, where $E\ge\Delta E$, it is at most
$e^{-q/2}$, with equality on an eigenvector of energy exactly $\Delta E$.
Hence
\begin{align}
    \|P_q-P_0\| & \le e^{-q/2},
    \label{eq:hastings07_gaussian_average_corrected}
\end{align}
and this is sharp when $\Delta E$ is attained. Since $\xi\ge2v/\Delta E$, the
choice of $q$ gives $q=l\Delta E/(6v)\ge l/(3\xi)$, so
(\ref{eq:hastings07_gaussian_average_corrected}) yields
$\|P_q-P_0\|\le e^{-l/6\xi}$. The exponent $l/3\xi$ printed in
(\ref{eq:hastings07_gaussian_average_source}) is twice the one that follows,
and with $\xi=2v/\Delta E$ and $\Delta E$ attained it is false.

\section{The corrected statement}

For every real $q$, the bound (\ref{eq:hastings07_gaussian_average_corrected})
holds, with the Gaussian of variance $\max(q/\Delta E^2,0)$ in place of the
normalized density.\footnote{\leanid{TNLean.MPS.AreaLaw.Hastings.norm_gaussianTimeAverage_sub_groundProjector_le}
    in \doclink{TNLean/MPS/AreaLaw/Hastings/GaussianTimeAverage}.}
With the source's choice of $q$ it gives $O(e^{-l/6\xi})$.

\section{Why the conclusion is unaffected}

The rest of the appendix already carries the weaker exponent. The spectral
cut-offs $O_L,O_R$ cost $O(e^{-l/6\xi})$, and this term is added in the
estimate that follows, so its error is
$O(\Delta E^{-2}J^2\sqrt q\,e^{-l/3\xi}+e^{-l/6\xi})$. The next estimate is
$O(\Delta E^{-2}J^2\sqrt q\,e^{-l/6\xi})$, and the final bound on
$\|O_BO_LO_R-P_0\|$ is
$O(\sqrt{l/\xi}\,\xi^2e^{-l/6\xi})$.\footnote{\path{References/0705.2024/law16.tex},
    lines~838--839, 875, 889 and~904--906.}
With the corrected exponent, the display that follows
(\ref{eq:hastings07_gaussian_average_source}) gains the additive term
$e^{-l/6\xi}$, which the estimate at line~875 already contains. No later
exponent changes, and Lemma~1 holds with the same $\xi'=6\xi$.

\bibliographystyle{alpha}
\bibliography{references}

\end{document}
